% Sviluppato con Texifier per macOS.
% Se si utilizza il suo generatore automatico LaTeX TexpadTeX prestare attenzione perch� alcuni pacchetti non funzionano correttamente, ad esempio hyperref.
% Una volta generato, se su Texifier, se si sta usando il generatore automatico LaTeX TexpadTeX, premere un tasto nel codice per generare l'indice, che altrimenti talvolta non appare automaticamente anche senza errori.

% Carattere dimensione 12.
\documentclass[12pt]{report}

% Per la stampa fronte-retro sostituire con:
% \documentclass[12pt, twoside]{report}

% Margini (4cm a sx, 2.5cm a dx, 2.5cm in alto, 2.5cm in basso).
\usepackage[top=2.5cm, bottom=2.5cm, left=4cm, right=2.5cm, centering]{geometry}

% Per la stampa fronte-retro sostituire con: 
% \usepackage[top=2.5cm, bottom=2.5cm, inner=4cm, outer=4cm, right=2.5cm, centering]{geometry}

% Interlinea.
\linespread{1.5}

% Librerie utili.
% Applicazione regole di scrittura per la lingua italiana.
\usepackage[italian]{babel}
% Codifica UTF-8. 
\usepackage[utf8]{inputenc}






% Matematica (matrici).
\usepackage{amsmath} 
% Links.
\usepackage{hyperref}
\hypersetup{   
    pdftitle={Relazione Progetto Esame Corso Laboratorio II - Giulio Nisi},
    pdfpagemode=FullScreen,
}
% Citazione.
\usepackage{epigraph}
% Grafo dipendenze files.
\usepackage{tikz} 
% Colora link GitHub.
\usepackage{xcolor}






% Stili pagina per il frontespizio e numerazione pagine.
\usepackage{scrlayer-scrpage} 
\ifoot[]{}
\cfoot[]{}
\ofoot[\pagemark]{\pagemark}
\pagestyle{scrplain}
% Font Times New Roman (simile).
\usepackage{mathptmx} 
% Per la formattazione dei titoli delle sezioni, capitoli, etc.
\usepackage{titlesec} 
% Formato delle intestazioni.
\titleformat{\chapter}[block]{\normalfont\LARGE\bfseries}{\thechapter.}{0.5em}{\LARGE}
\titlespacing*{\chapter}{0pt}{-20pt}{25pt}

\begin{document}

% Frontespizio.
\begin{titlepage}
\begin{figure}
    \centering\includegraphics[scale=0.5]{cherubino-logo-unipi.png}
\end{figure}

\begin{center}
    {\LARGE{ Corso di Laurea in Informatica \\}}
    \vspace{2cm}
    {\Large { PROGETTO ESAME CORSO LABORATORIO II }}\\
    \vspace{2cm}
    {\Large { "Il Paroliere" videogioco online scritto in C, multi-threaded, CLI, POSIX, client/server }}
\end{center}

\vspace{2cm} 

\begin{minipage}[t]{0.47\textwidth}
	{\large{\bf Professori:\\ Patrizio Dazi, Luca Ferrucci}}
	\vspace{0.5cm}

\end{minipage}\begin{minipage}[t]{0.47\textwidth}\raggedleft
	{\large{\bf Candidato: \\ Giulio Nisi\\ }}
\end{minipage}

\vspace{25mm}

\centering{\large{\bf ANNO ACCADEMICO 2023/2024 }}
\end{titlepage}
% Fine frontespizio.

% Indice e numerazione pagine.
\tableofcontents
\thispagestyle{empty}
\clearpage
\setcounter{page}{1} 

% Inizio contenuto documento.
% Note.
% Per nasconderlo dalla numerazionie nell'indice.
\addcontentsline{toc}{chapter}{Note} 
% File note.
% Asterisco per nascondere numerazione.
\chapter*{Note}

Si consiglia la lettura del codice in caso di necessit\'a, \'e stato molto commentato, forse pi\'u del necessario e potrebbe risultare pi\'u chiaro di tale relazione riassuntiva. Il codice \'e scritto in inglese per essere compreso potenzialmente da un pubblico pi\'u ampio e pubblicato su {\color{cyan}\href{https://github.com/JuliusNixi/Boggle}{GitHub}}. Per ottemperanza alle richieste, le specifiche scritte nel testo del progetto sono tutte soddisfatte anche in italiano, ad esempio, il client accetta anche i comandi in italiano (oltre a quelli in inglese), ed entrambe gli eseguibili gestiscono gli argomenti passati da riga di comando italiani (oltre a quelli in inglese). Il codice \'e stato sviluppato su macOS (Sonoma 14.5) (Darwin Kernel Version 23.5.0) (arm64) (sfruttando le chiamate POISX) e testato sulla macchina di laboratorio Linux (laboratorio2.di.unipi.it).

\vspace{5mm}

\epigraph{"Scrivete programmi che facciano una cosa e che la facciano bene. Scrivete programmi che funzionino insieme. Scrivete programmi che gestiscano flussi di testo, perch\'e quella \'e un'interfaccia universale."}{\textit{Peter H. Salus, A Quarter Century of Unix, Unix philosophy. }}



\chapter{Strutture dati ed algoritmi}
\label{ch:strutturedatiedalgoritmi}

% --- Inizio del Capitolo 1 ---

\iffalse

Esempio di una nota\footnote{CleanCode} con citazione a sezione~\ref{sec:server}.

\begin{figure}[H]
    \centering
    \includegraphics[width=0.2\textwidth]{immagini/logo-dip_blu_hr.png}
    \caption{Esempio di un'immagine}
    \label{fig:immagine1}
\end{figure}

\begin{lstlisting}[caption={Esempio codice JavaScript},label={Esempio codice JavaScript}, language=JavaScript]
function foo(num) {
    const bar = 2;
    return num + bar;
}

const result = foo(2);
// result -> 4
\end{lstlisting}


\fi

\section{Server}
\label{sec:serverstrutturedatiedalgoritmi}

% Sezione 1... con citazione bibliografica~\cite{CleanCode}

\subsection{Lista giocatori}
\label{subsec:listagiocatori}

La struttura dati principale usata per la gestione dei clients/giocatori, \'e una lista concatenata di elementi con una struttura e puntatore al prossimo elemento della lista. \'E stata scelta per:
 \\ - Possibilit\'a di gestire potenzialmente infiniti giocatori, risorse computazionali permettendo.
 \\ - Allocazione dinamica on demand che permette di sprecare poca memoria (non preallocandola) quando i giocatori si scollegano (la memoria viene liberata), o sono assenti.
 \\ - Il vantaggio che tutte le operazioni riguardanti un solo determinato client possono essere svolte dal relativo thread attraverso il puntatore alla struttura rappresentante suddetto giocatore e qui si hanno tutte le informazioni che potrebbero servire.
  \\
 Lo svantaggio sostanziale \'e che talvolta dobbiamo scorrere tutta la lista sincronizzandoci opportunamente con particolare attenzione, perch\'e, in caso contrario, a differenza di un array per esempio, potrebbero verificarsi problemi molto pi\'u gravi in certe fasi, come nell'aggiunta o rimozioni di elementi. Per far ci\'o si utilizza (ma non \'e l'unico) listmutex.
 \\
 Tale lista \'e condivisa tra clients registrati e non.
 \\
 
 
 \subsection{Matrice di gioco}
 \label{subsec:matricedigioco}

La matrice di gioco, avendo dimensione predeterminata alla compilazione e costante \'e allocata staticamente globalmente come matrice di caratteri (char[][]). Il 'Qu' viene salvato come 'Q', ma gestito opportunamente.
\\
\subsection{Parole e dizionario}
\label{subsec:paroledizionario}
Le parole del file dizionario sono memorizzate in un array (words) di stringhe (char**) allocato dinamicamente, sia l'array, che le stringhe stesse. Le parole sono caricate tutte in memoria, leggendo il file, solo una volta inizialmente, limitando l'impatto delle costose operazioni di I/O con il disco. Vi \'e poi una copia di words (words\_valid) allocata in memoria, dinamicamente, per\'o, NON riallocando le stringhe ed utilizzando inutilmente altra memoria, ma semplicemente copiando i puntatori da words. Ad ogni inizio nuovo gioco, contestualmente al cambio della matrice di gioco, i puntatori alle stringhe di words\_valid vengono aggiornati, quelli delle parole presenti nella nuova matrice, rimangono invariati (copiati da words), mentre quelli delle parole NON presenti, vengono incrementati fino a raggiungere il terminatore (\textbackslash0) di stringa. Vi \'e poi un'altra copia di words\_valid (words\_validated), allocata identicamente, per ogni giocatore registrato, il cui puntatore si trova nella struttura rappresentate tale giocatore. Quando una parola valida (presente in words\_valid E in words\_validated) viene sottomessa da un utente registrato, per la PRIMA volta, words\_validated viene aggiornato, incrementandone il puntatore fino al terminatore, cos\'icch\'e ad nuova sua sottomissione della medesima parola, essa possa essere rifiutata senza attribuirne nuovamente il punteggio al player.
\\
Riassumendo, le parole del file dizionario vengono tutte cercate ad ogni cambio di matrice di gioco (inizio nuovo gioco) (ma NON lette dal file, vengono caricate tutte solo la prima volta in words) nella matrice stessa per aggiornare words\_validated, del quale ogni giocatore (registrato) avr\'a una copia personale per tracciare quelle gi\'a sottomesse. 
\\
\subsection{Algoritmo di ricerca}
\label{subsec:algoritmoricerca}
Per gestire la ricerca delle parole proposte si \'e scelto di utilizzare il seguente algoritmo. Per ogni parola viene visitata tutta la matrice e per ogni suo carattere (della matrice) invocata la funzione searchWordInMatrix(), la quale parte dal carattere della matrice specificato (con indici i di riga e j di colonna della corrente iterazione) e ricorsivamente controlla che i caratteri adiacenti [(i, j+1),(i+1, j),(i,j-1),(i-1,j)] nella matrice siano corrispondenti al prossimo carattere di parola cercato. Se non vi \'e corrispondenza, o i o j sono fuori dal range di grandezza della matrice, si ritorna un fallimento, altrimenti, se il prossimo carattere di parola voluto \'e identico al terminatore di stringa, allora significa che precedentemente abbiamo trovato tutte le lettere disposte adiacentemente e quindi ritorniamo un successo. Bisogna prestare attenzione a "marcare" i caratteri gi\'a visitanti per evitare di poter utilizzare pi\'u volte lo stesso carattere della matrice nella composizione di una parola o di perdersi in loop infiniti. Per quanto riguarda la complessit\'a? Vediamo un esempio. Per semplicit\'a assumiamo di avere M una matrice N x N (utilizzando sempre matrici quadrate), di avere K parole sul file dizionario, e di avere una parola W, da cercare, in input, ossia verificarne la presenza nella lista delle parole E nella matrice. Con la predetta soluzione, dopo aver letto dal file le parole ed averle inserite in words, si dovr\'a riempire words\_valid (solo una volta ad ogni inizio gioco), scorreremo quindi words in K passi, ad ogni passo iteriamo su M effettuando $N \cdot N$ passi, invocando ad ognuno, searchWordInMatrix(), la quale per trovare la parola al massimo effettuer\'a proprio altri $N \cdot N$ passi che rappresentano la parola di lunghezza massima che si pu\'o trovare nella matrice, vediamo a seguire un esempio con N = 5 dove con la numerazione \'e riportato solamente uno dei possibili percorsi completi che searchWordInMatrix() potr\'a intraprendere per trovare\\\\ la parola di lunghezza massima 25:
$
\begin{bmatrix}
1 & 2 & 3 & 4 & 5\\
10 & 9 & 8 & 7 & 6\\
11 & 12 & 13 & 14 & 15\\
20 & 19 & 18 & 17 & 16\\
21 & 22 & 23 & 24 & 25\\
\end{bmatrix}
$.\\\\
Quindi ricapitolando $ K \cdot N \cdot N \cdot N \cdot N = N^4K $ la pima volta, a seguire, ogni parola cercata da un giocatore comporter\'a solamente la ricerca in words\_valid (ed il controllo in words\_validated, ma questo si far\'a ad accesso diretto se trovata, dato che gli array sono allineati) con costo $ K $. In conclusione credo sia una discreta implementazione, che sfrutta la potenza dell'aritmetica dei puntatori, le stringhe sono effettivamente in memoria solo una volta e a fronte di un costo iniziale pi\'u elevato permette di rispondere pi\'u velocemente a seguire ai giocatori.
\\
\subsection{Coda di fine gioco}
\label{subsec:coda}
Per la struttura del progetto vista, la coda richiesta dalle specifiche, non sarebbe stata necessaria, anzi \'e risultata un'inutile complicazione. Sarebbe stato sufficiente sbloccare i threads dei clients e segnalarli di inviare il messaggio di fine gioco con la scoreboard al client gestito. Ogni thread cos\'i, sfruttando la peculiaret\'a sopra spiegata, avrebbe potuto inviare il messaggio utilizzando il puntatore alla struttura del giocatore in possesso che fornisce gi\'a tutte le informazioni necessarie. Comunque per aderenza alle richieste, la coda \'e stata utilizzata. Essa \'e stata implementata come lista concatenata, ogni elemento \'e una struttura contenente un puntatore al corrispettivo client, un puntatore ad un messaggio (con struttura come richiesta) contenente nel campo data, il nome del client (se registrato, un placeholder altrimenti) ed il punteggio ottenuto nel gioco separati con una virgola, ed infine un puntatore all'elemento successivo. Per la sincronizzazione viene utilizzato (ma non \'e l'unico) queuemutex.
	
\section{Client}
\label{sec:clientstrutturedatiedalgoritmi}
Nel client vengono utilizzate due liste allocate dinamicamente. La prima di una struttura contenente stringhe allocate dinamicamente, dove ogni stringa ha una lunghezza massima BUFFER\_SIZE. La seconda di una struttura contenente messaggi che rappresentano le risposte ricevute dal server. Per sapere il peculiare motivo di queste fare riferimento alle pagine seguenti. Nessun particolare algoritmo \'e utilizzato.


























\chapter{Struttura codice}

Il codice si trova tutto nella cartella "./Src/" (soruces). Suddiviso in tre sottocartelle "./Server", "./Client" e "./Common". Le dipendenze sono le seguenti:
\\
\\
"boggle\_server" ("paroliere\_srv") eseguibile dipende da "boggle\_server.c", "server.c", "common.c".
\\
"boggle\_client" ("paroliere\_cl") eseguibile dipende da "boggle\_client.c", "client.c", "common.c".
\\
\\
"boggle\_server.c", "server.c" dipendono da "server.h".
\\
"boggle\_client.c", "client.c" dipendono da "client.h".
\\
\\
"server.h", "client.h" dipendono da "common.h".
\\
\\
"common.c" dipende da "common.h".
\\
\\
Si noti che "common.h" e "common.c" includono librerie, dichiarazioni ed implementazioni comuni a server e client. "boggle\_server.c" e "boggle\_client.c" contengono solamente i setup e lo scheletro del loro funzionamento. L'effettiva complessa e sostanziale implementazione \'e demandata rispettivamente ai "server.c" e "client.c" files. "server.h" contiene le dichiarazioni di funzioni e le strutture dati condivise tra "boggle\_server.c" e "server.c". Similmente per il client.
\\
\\
A seguire un diagramma (o grafo) delle dipendenze per il server. Quello del client \'e analogo.


% Grafo dipendenze file.
\usetikzlibrary{shapes,arrows,positioning}

% Setup.
\begin{tikzpicture}[
    node distance = 2cm and 3cm,
    file/.style = {rectangle, draw, rounded corners, minimum width=3cm, minimum height=2cm, align=center},
    arrow/.style = {->, thick}
]

% Primo livello.
\node[file] (boggle_server) {boggle\_server};

% Secondo livello.
\node[file, below left=of boggle_server] (boggle_server_c) {boggle\_server.c};
\node[file, below=of boggle_server] (server_c) {server.c};
\node[file, below right=of boggle_server] (common_c) {common.c};

% Terzo livello.
\node[file, below=of server_c] (server_h) {server.h};

% Quarto livello.
\node[file, below=of server_h] (common_h) {common.h};

% Connessioni.
\draw[arrow] (boggle_server) -- (boggle_server_c);
\draw[arrow] (boggle_server) -- (server_c);
\draw[arrow] (boggle_server) -- (common_c);
\draw[arrow] (server_c) -- (server_h);
\draw[arrow] (boggle_server_c) -- (server_h);
\draw[arrow] (server_h) -- (common_h);
\draw[arrow] (common_c) -- (common_h);

\end{tikzpicture}
% Fine grafo dipendenze file.
\leavevmode
I files eseguibili si trovano nella cartella "./Bin" e sono compilati qua, alla fine copiati da make nella root del progetto con alias "./paroliere\_srv" e "./paroliere\_cl".
\\
I files dizionario si trovano in "./Data/Dicts/".
\\
I files matrici si trovano in "./Data/Matrices/".
\\
 Tutto il materiale di tests and debug si trova in "./Tests/".
\\
Vi \'e una cartella "./Studies/". Contiene degli snippets di codice per fare prove e tests usati per assicurarsi del funzionamento di alcune meccaniche prima di implementarle nel progetto con il rischio di introdurre tediosi bugs.





\chapter{Struttura programmi}

\section{Server e client in comune}

\subsection{Funzioni comuni}

Il server ed il client condividono librerie utilizzate e molto codice per delle funzionalit\'a richieste comuni grazie ai files common.h e common.c. 
\\
Enunciamo sinteticamente le funzioni condivise. Abbiamo una funzione per fare il parsing dell'IP ottenuto dalla riga di comando, una per trasformare le stringhe in UPPERCASE o lowercase, le fondamentali receiveMessage() e sendMessage() che servono a scambiarsi i messaggi tra server e client, una funzione destroyMessage() per deallocare tali messaggi, due wrapper mLock() e mULock(), rispettivamente delle funzioni pthread\_mutex\_lock() e pthread\_mutex\_unlock(), utilizzati per controllare gli errori, una funzione handleError() dedicata alla gestione degli errori, una funzione wrapper printff() di quella della libreria, printf(), per stampare messaggi su stdout ed errori su stderr, che insieme ad un mutexprint permette la stampa di stringhe multiple senza che ne vengano interposte altre da altri thread quando necessario, una funzione createBanner() che crea delle stringhe carine da stampare per suddividere l'output e facilitarne la lettura ed infine una funzione threadSetup() che si occupa di associare un distruttore ad ogni thread, richiamato alla sua terminazione. Tali funzioni sono identiche in server e client.
 
 \subsection{Funzioni comuni ma con differente implementazione}

 I thread condividono il possedere un thread dedicato alla gestione dei segnali ricevuti. L'implementazione del thread differisce leggermente tra i due, poich\'e il server aveva la necessit\'a di gestire il SIGALRM in modo impegnativo (perch\'e usato per terminare il gioco), che il client non doveva. I server ed il client hanno due versioni differenti della funzione threadDestructor() che come predetto viene chiamata alla terminazione di qualunque thread, ed in base al chiamante esegue dei cleaner differenti. Per concludere vi \'e una funzione atExit() che viene chiamata al termine del programma dal thread main che si occupa anch'essa di fare pulizia.
 
\section{Gestione segnali}

Come anticipato, vi \'e un thread dedicato alla gestione dei segnali. Questo intercetta i segnali ricevuti da gestire, attraverso un loop, con una sigwait(). Inizialmente lo sviluppo \'e stato fatto con la registrazione dei signals handlers, ma questo esponeva ad innumerevoli difficolt\'a dovute al fatto che il segnale pu\'o finire ad un thread casuale e sopratutto che nel suo gestore, si \'e costretti ad utilizzare un numero estremamente ridotto di funzioni che devono avere la caratteristica di essere async-signal-safe. Con l'uso di un thread dedicato e della sigwait() invece sappiamo gi\'a il thread che ricever\'a il segnale e possiamo gestirlo liberamente.

\section{Server}

\subsection{Thread main}

Il server, dopo un'iniziale fase di setup, con inizializzazioni di varie strutture dati, registrazione del thread distruttore e della funzione atExit(), controllo degli argomenti ricevuti da riga di comando, avvio e configurazione del thread per la gestione dei segnali, registrazione del gestore dell'unico segnale (SIGUSR1) non gestito dal predetto thread, caricamento del file dizionario e matrici (se presente), apre un socket, si mette in ascolto, avvia il primo gioco (impostando un timer di fine partita con la chiamata alarm()) e si mette ad accettare clients indefinitivamente in un ciclo. All'arrivo di un client (connessione), lo accetta, inizializza un nuovo elemento della lista clients, aggiungendolo a tale lista (con dovute sincronizzazioni) ed infine avvia un thread (clientHandler()) dedicato alla gestione del giocatore, passandogli come argomento il puntatore alla propria struttura dati aggiunta alla lista.

 \subsection{Thread clientHandler()}

 Il neoavviato thread si mette in un loop ad attendere la ricezione di un messaggio dal client attraverso la funzione receiveMessage() che al suo interno utilizza la chiamata read() per leggere le varie componenti del messaggio dal socket utente. All'arrivo di un messaggio COMPLETO tenta di acquisire il proprio mutex. Ogni client, infatti, all'interno della propria struttura dati, ha anche un suo personale mutex, che acquisisce durante l'elaborazione di un messaggio ricevuto, per poi rilasciarlo al termine. Tale mutex \'e posto DOPO la lettura del messaggio, ma PRIMA della sua elaborazione.
 
 \subsection{Thread signalsThread()}

Quando il tempo di gioco termina, si riceve un segnale SIGALRM che verr\'a intercettato dal thread gestore segnali ed inizier\'a la fase pi\'u complessa del progetto. Sincronizzandosi opportunamente bloccher\'a tutti i threads clientHandler() acquisendo ciascun loro mutex. Se non dovesse riuscire ad acquisirne uno, ritenter\'a. In questo modo si \'e scelto di dare la priorit\'a al thread clientHandler() che sta gestendo un messaggio, questo ci garantisce la piacevole propriet\'a che tutte le richieste ricevute PRIMA dello scadere del timer siano processate come \'e giusto che sia. A seguire, il thread signalsThread() abilita la pausa, modificando una variabile globale (pauseon), che in seguito istruir\'a i clientHandler() threads di rispondere alle richieste di conseguenza (siamo in pausa). Ad ognuno di questi threads da parte del signalsHandler() che sta gestendo il SIGALRM verr\'a inviato il segnale SIGUSR1 per il quale nel main era stato registrato un handler (questo segnale NON \'e gestito dal signalThread(), ma ogni clientHandler() thread lo ricever\'a e gestir\'a). Lo scopo di questa azione \'e interrompere eventuali read() sulle quali il clientHandler() si dovesse trovare, perch\'e adesso (dopo esser stati sbloccati dal signalsHandler() con il rilascio del proprio mutex) i clientHandler() threads dovranno riempire la coda con i messaggi contenenti il punteggio del giocatore con le dovute sincronizzazioni (usando il queuemutex). Quando la coda sar\'a stata riempita, il signalsThread() thread avvier\'a il thread scorer() che stiler\'a la classifica finale recuperando i messaggi dalla coda ed ordinandoli per punteggio discendente con una chiamata qsort(). Al termine del scorer() thread, il signalsThread() che lo avr\'a aspettato con una join, ribloccher\'a tutti i clientHandler() threads (sempre acquisendo ciascun loro mutex scorrendo la lista), li comunicher\'a che la scoreboard CSV \'e pronta e loro (i clientHandler()) dopo esser stati tutti liberati, invieranno la scoreboard ai propri (corrispettivi) utenti. A questo punto signalsThread() non dovr\'a far altro che avviare il thread gamePause() che eseguir\'a la sleep di durata della pausa, signalsThread() far\'a il join su questo thread. Nel frattempo ovviamente tutti i clientHandler() saranno liberi di operare rispondendo liberamente alle richieste dei giocatori (tenendo conto che la pausa \'e stata abilitata). Al termine di gamePause() thread, signalsThread() bloccher\' nuovamente tutti i clientHandler() threads, disabiliter\'a la pausa, imposter\'a una nuova matrice (dal file o causale), aggiorner\'a il words\_validated array cercando tutte le parole del dizionario nella nuova matrice, imposter\'a di conseguenza il words\_validated di ciascun giocatore ed infine avvier\'a un nuovo gioco (nuova chiamata ad alarm()) e si metter\'a in attesa di nuovi segnali, cos\'i il ciclo potr\'a ricominciare.

\section{Client}

Lo sviluppo del client \'e stato sorprendentemente pi\'u arduo di quello del server, nonostante la sua abbastanza oggettiva semplicit\'a. Il concetto chiave su cui mi sono concentrato \'e stato la pulizia della stampa delle risposte del server e la successiva stampa (e gestione degli input) del prompt. Pu\'o sembrare superficiale, ma senza la giusta attenzione, l'interfaccia grafica risultava disastrosa. Pu\'0 sembrare un dettaglio, ma \'e stato difficile far fronte a questa questa esigenza,  sulla quale mi sono incaponito nel non voler scendere a compromessi. Nello specifico, il problema e\' che quando viene ricevuta e stampata una risposta dal server, l'utente potrebbe essere nel mezzo di una digitazione e aver gi\'a inserito dei caratteri nel buffer STDIN. Per vedere il problema si pu\'o guardare il codice nel file brokeninputoutput.c, nella cartella Studies/. Ho individuato due modi per affrontare il problema:
\\
- Prevenirlo. Sviluppando il client come un singolo thread, che attende che l'utente completi l'input (premendo ENTER, che tramite '\textbackslash n' interrompe la read()), poi esegue l'input ed infine controlla se ci sono risposte del server da stampare e ricomincia il ciclo. Oppure, sviluppare il client come multithread e sincronizzare il thread che gestisce l'input con quello che stampa le risposte del server. Entrambe le soluzioni, hanno lo stesso difetto: la read() deve essere obbligatoriamente bloccante, per essere sicuri di leggere l'intero buffer STDIN.
\\
- Risolverlo, lasciando che il buffer STDIN possa "sporcarsi" e pulirlo quando necessario. Sembra semplice, ma non lo \'e, ho trovato solo metodi non funzionanti su internet.
\\
Ho finito per tentare con una nota libreria chiamata ncurses. Questa permette di gestire il terminale nel dettaglio, ma era inappropiata. Risolveva il problema, ma generava una complessit\'a insostenibile per le semplici stampe che dovevo fare. 
\\
Dopo innumerevoli tentativi sono riuscito a risolvere il problema in modo semplice, ricorrendo a due librerie ma senza doverne fare un uso estensivo. Ho utilizzato termios, vedere termiostest.h nella cartella Studies/ e fcntl. Utilizzo due threads, uno per la gestione dell'input utente e della stampa delle risposte del server (quest'ultime vengono lette da una lista) ed uno che legge le risposte del server dal socket e le inserisce in lista. Interrompo la read() nel primo thread molto spesso e controllo se ci sono risposte del server, in caso affermativo, le stampo, ripulisco lo STDIN con una getchar() NON BLOCCANTE e ritorno al prompt pulito, in modo che se l'utente avesse digitato qualcosa di incompleto possa riscriverlo e/o modificarlo (a seguito della risposta). Invece, se non ci sono risposte dal server, la read() riprende con lo STDIN inserito dall'utente che non si accorge di nulla. L'unico svantaggio \'e che le stampe delle risposte del server non avvengono in modo asincrono, si ottiene un effetto simile per l'utente, perch\'e la read() viene interrotta ogni pochi millisecondi per stampare le risposte del server. La ricezione delle risposte e l'inserimento in lista invece e\' eseguita da un thread ad-hoc e quindi e\' totalmente asincrono, ma la stampa no. La sincronizzazione tra i due threads avviene con un mutex. Le due librerie citate servono solo a rendere la getchar() non bloccante e permettermi quindi di pulire l'STDIN. Per approfondire vedasi inputoutputasynctests.c nella cartella Studies/.
\\
Un ulteriore problema affrontato \'e stato quello di non sapere anticipatamente la grandezza dell'input dell'utente. Non vi sono infatti limitazioni sulla lunghezza della parola o del nome utente, per poter gestire input di arbitraria dimensione si leggono BUFFER\_SIZE caratteri dallo STDIN con una read(), essi vengono inseriti in un char array[BUFFER\_SIZE] statico, ma da qui vengono poi copiati in una stringa allocata sullo heap dinamicamente (sempre di lunghezza BUFFER\_SIZE) che viene aggiunta ad una lista concatenata di stringhe. Scegliendo cos\'i un BUFFER\_SIZE adeguato allocheremo sempre la dimensione di memoria che pi\'u si avvicina alla grandezza dell'input dell'utente evitando di preallocare tanta memoria inutilmente che poi comunque se allocata staticamente potrebbe terminare (in caso di input utente grandi), con una lista di stringhe si ovvia a questo problema.




\chapter{Struttura tests}

Inizialmente l'idea era di scrivere degli unit tests. Nello specifico cercare di rendere il progetto il pi\'u modulare possibile attraverso la scrittura di funzioni separate e testarne input ed output man a mano che si progrediva nello sviluppo. Ho intrapreso questa strada, testimone il fatto che nella cartella "./Tests/C/" ho lasciato il file "tests.c", il quale funziona, si pu\'o compilare ed eseguire. Tuttavia, con l'avanzare dello sviluppo, mi sono reso conto che per come ho strutturato il progetto era disastroso. Avrei potuto testare poche singole funzioni ma non ad un livello soddisfacente nell'insieme. Complice il fatto di avere i codici divisi in pi\'u file ed di aver utilizzato per molte funzionalit\'a variabili globali condivise. Allora, ho guardato qualche framework di testing in C, ma non ho trovato nulla che facesse al caso mio. Alla fine, ho pensato di fare i tests in modo esterno, pensando che forse questo approccio era proprio quello voluto per sondare altri argomenti trattati nel corso. Ho cos\'i creato un programma distinto che creasse tanti processi client, simulando proprio il comportamento di utenti ed interazioni con il server su larga scala. Ha funzionato benissimo! Ad ogni modo credo sia stato didatticamente utile perch\'e ho dovuto maneggiare descrittori, files, pipes e segnali. In questo modo ho potuto eseguire moltissimi tests che sono stati vitali e mi hanno permesso di correggere innumerevoli bugs che con pochi clients (testati manualmente) non si verificavano, soprattutto sviste con la memoria e la sincronizzazione. Per eseguire questi tests \'e sufficiente:

\begin{enumerate}
  \item Avviare il server con i parametri desiderati.
  \item Dalla root del progetto, dopo aver compilato con "make", eseguire "./paroliere\_tests server\_ip porta\_server". Oppure "make tests" per utilizzare IP e porta default.
\end{enumerate}
\leavevmode
Ci sono alcuni parametri all'inizio del file che possono essere modificati se voluto. Il file crea tantissimi processi client e per ognuno di essi svolge un numero di azioni che possono essere esattamente tutte quelle che un vero utente potrebbe eseguire. Inviare comandi non validi, registrarsi, richiedere la matrice, sottomettere una parola, di quelle valide e non valide, uscire e persino disconnettersi forzatamente. Tutto viene svolto casualmente con un intensivo uso di random (ma comunque replicabile tramite seed). Tra ogni azione si aspetta un numero casuale di millisecondi per aumentare l'entropia nelle sincronizzazioni. I comandi vengono sottomessi, creando il processo client con una chiamata "fork()", collegandosi tramite pipe al suo STDIN. L'input inviato (comandi spediti al server) e l'output ricevuto (le risposte che il client ha ricevuto dal server) vengono salvati in dei file di logs in "./Tests/C/Logs/" per potere esaminare che tutto funzioni correttamente.



\chapter{Compilazione ed esecuzione}

Il codice pu\'o essere compilato sia nella root del progetto, sia nella cartella Src/.
\\
\\
In entrambi i casi valgono i seguenti comandi:
\\
- make = Compila tutto: client, server and tests.
\\
- make execs = Compila ed esegue il server, con gli argomenti di default (si veda il Makefile in Src/).
\\
- make execc = Compila ed esegue il client, con gli argomenti di default (si veda il Makefile in Src/).
\\
- make clean = Rimuove tutti i file oggetto, eseguibili, file di logs dei tests e termina tutti gli eventuali processi in esecuzione di server e client.
\\
- make tests = Compila il server e compila ed esegue il file C tests.
\\
\\
L'utilizzo degli eseguibili si pu\'o svolgere facendo riferimento al seguente formato:
\\
Server Usage: boggle\_server server\_ip server\_port [--matrices matrices\_filepath] [--duration game\_duration\_in\_minutes] [--seed rnd\_seed] [--dic dictionary\_filepath].
\\
Client Usage: boggle\_client server\_ip server\_port.
\\
\\
\textbf{ATTENZIONE:} L'utilizzo di alcuni percorsi assoluti hard-coded (ad esempio il percorso del file dizionario di default quando non specificato da CLI) assume che la \textbf{WORKING DIRECTORY} sia  \textbf{BIN} causa la presenza nel percorso di "../". Utilizzando make nella root non si incorre mai in problemi perch\'e si occupa lui di cambiare directory quando esegue, ma se si avvia i binari manualmente bisogna tenere conto di questo. TL;DR: Fare sempre un cd in Bin/ prima di eseguire.



\chapter{Miscellanea}

Le funzionalit\'a di cui sono pi\'u orgoglioso sono:
\begin{itemize}
\item Il server pu\'o far giocare potenzialmente infiniti giocatori.
\item La matrice di gioco pu\'o essere N x N di dimensione arbitraria.
\item Le parole del gioco possono avere lunghezza arbitraria.
\item L'input dell'utente nel client pu\'o essere di lunghezza arbitraria (parole sottomesse e lunghezza nickname di registrazione).
\item Le stampe nel client sono ben strutturate, la GUI non mostra mai stampe incoerenti.
\item I tests eseguiti mostrano una buona resa e stabilit\'a di server e clients.
\item Il client \'e semplice poich\'e il grosso del lavoro lo esegue il server.
\end{itemize}
\leavevmode 
Tutta la sincronizzazione \'e stata volutamente fatta solamente tramite mutex, dove talvolta le variabili di condizione o i semafori si sarebbero potuti utilizzare, per cercare di rendere il codice pi\'u semplice ed intuitivo possibile, seppur abbia ancora le sue complessit\'a.
 \\
 Le parole, indipendentemente da come arrivino dal client e da come siano scritte nei file dizionario vengono convertite in UPPERCASE dal server. Anche i caratteri componenti la matrice di gioco sono gestiti esclusivamente in UPPERCASE, sia nella generazione casuale, sia nella lettura da file (se scritte in lowercase in quest'ultimo vengono convertite in UPPERCASE). Gli unici caratteri ammessi per le parole, le matrici e i nicknames degli utenti sono: "abcdefghijklmnopqrstuvwxyz", contenuto nella definizione "ALPHABET" in "server.h". L'input del client, invece, viene tutto convertito a lowercase prima di essere elaborato ed inviato, quindi il comando inserito "MaTrIcE" sar\'a valido come "matrice".
 \\
 A tutte le informazioni che possono essere espresse solamente con numeri interi positivi \'e stato assegnato il tipo "unsigned long int" abbreviato con il "typedef" ad "uli". Per le variabili booleane (o a carattere discreto) \'e stato utilizzato il tipo "char" poich\'e occupa solo un byte.
 \\
 Perch\'e nel server le "read()" nella funzione "receiveMessage()", chiamata dai "clientHandler()" threads, viene interrotta tramite l'invio (da parte del "signalsThread()" thread) del segnale SIGUSR1, mentre nel client la "read()" dello STDIN viene interrotta autonomamente grazie alla libreria "fcntl"? Semplicemente perch\'e non sapevo si potesse fare con questa libreria quando ho sviluppato il server, che ho sviluppato per prima, quando successivamente, sviluppando il client l'ho scoperto, sono rimasto deluso da tutta l'inutile complessit\'a che ho dovuto gestire nel server, ma questo funziona ugualmente e mostra pi\'u attinenza agli argomenti trattati nel corso, quindi l'ho felicemente lasciato cos\'i.
 \\
 Durante lo sviluppo del client, per tentare di risolvere il problema spiegato sulla sincronizzazione input/output e buffer STDIN, sono finito per sperimentare con una nota libreria chiamata "ncurses". Questa permette di gestire il terminale nel dettaglio, ma era inapropiata al mio use case. Risolveva il problema, ma generava una complessit\'a in tutto il resto insostenibile per le semplici stampe che dovevo fare. Questo mi ha insegnato che anche fare delle "presumibilmente semplici" interfacce grafiche testuali non \'e sempre cos\'i scontato come credessi.

 
\iffalse
\begin{figure}[H]
    \centering
    \includegraphics[scale=0.7]{immagine.png}
    \caption{Descrizione.}
\end{figure}
\fi

 

\end{document}

